\documentclass[11pt, a4paper]{../problems_template}

\begin{document}

\begin{titlepage}
  \begin{center}
	\Huge{Журнал Квант}
  \end{center}
\end{titlepage} 

\chapter{1970}

\begin{exercise}[7]
$a$, $b$, $c$ - длины сторон треугольника. Докажите неравенство
$$
\frac{a}{b + c - a} + \frac{b}{c + a - b} + \frac{c}{a + b - c} \geqslant 3
$$
\end{exercise}

\begin{Solution}
Тогда $a = x + y$, $b = y + z$, $c = z + x$.
$$
\frac{a}{b + c - a} + \frac{b}{c + a - b} + \frac{c}{a + b - c} = \frac{x + y}{2z} + \frac{y + z}{2x} + \frac{z + x}{2y} = \frac{1}{2}\left( \frac{x}{z} + \frac{z}{x} \right) + \frac{1}{2} \left( \frac{y}{z} + \frac{z}{y} \right) + \frac{1}{2}\left( \frac{y}{x} + \frac{x}{y} \right) \geqslant 3
$$
\end{Solution}

\begin{exercise}[24]
Любую дробь $\frac{m}{n}$, где $m$ и $n$ - натуральные числа, $1 < m < n$, можно представить в виде суммы нескольких дробей вида $\frac{1}{q}$, причём таких, что знаменатель каждой следующей из этих дробей делится на знаменатель предыдущей дроби. Докажите это.
\end{exercise}

\begin{Solution}
По условию задачи
$$
\frac{m}{n} = \frac{1}{a_{1}} + \frac{1}{a_{1}a_{2}} + \frac{1}{a_{1}a_{2}a_{3}} + \dots + \frac{1}{a_{1}a_{2}\dots a_{k}} = \frac{1}{a_{1}} \left( 1 + \frac{1}{a_{2}} \left( 1 + \frac{1}{a_{3}} (1 + \dots ) \right) \right)
$$
Исходя из этого можно предложить следующий алгоритм. Пусть $a_{1}$ будет  минимальным числом при котором выполняется $ma_{1} > n$. Очевидно, что $ma_{1} - n < m$. Тогда 
$$
\frac{ma_{1}}{n} - 1 = \frac{ma_{1}-n}{n} = \frac{m_{1}}{n} < \frac{m}{n}
$$
Аналогично поступая с $\frac{m_{1}}{n}$ находим $a_{2}$. Так как числитель всегда уменьшается и взаимно прост со знаментателем, то в итоге он станет равным единице. На этом алгоритм завершается. 
\end{Solution}

\begin{exercise}[27]
Если
$$
\frac{a}{b-c} + \frac{b}{c-a} + \frac{c}{a-b} = 0
$$
то
$$
\frac{a}{(b-c)^{2}} + \frac{b}{(c-a)^{2}} + \frac{c}{(a-b)^{2}} = 0
$$
\end{exercise}

\begin{exercise}[M33]
Рассмотрим натуральное число $n > 1000$. Найдём остатки от деления числа $2^{n}$ на числа $1, 2, 3, \dots, n$ и найдём сумму всех этих остатков. Докажите, что эта сумма больше $2n$.
\end{exercise}

\begin{Solution}

\end{Solution}

\chapter{1971}

\begin{exercise}[63]
Можно ли из 18 плиток размером $1 \times 2$ выложить квадрат так, чтобы при этом не было ни одного прямого "шва", соeдиняющего противоположные стороны квадрата и идущего по краям плиток?
\end{exercise}

\begin{Solution}
Предположим, что таких "швов" нет. Тогда каждый из пяти вертикальных разрезов квадрата должен пересекать не менее двух горизонтальных плиток. Аналогично для горизонтальных разрезов. Значит всего должно быть не менее 10 вертикальных и 10 горизонтальных плиток. Но их всего 18. Противоречие. 
\end{Solution}

\chapter{1974}

\begin{exercise}[287]
Существует ли такая последовательность натуральных чисел, что любое натуральное число представимо в виде разности двух чисел этой последовательности единственным образом?
\end{exercise}

\begin{Solution}
Рассмотрим алгоритм для построения такой последовательности. Пусть последовательность уже построена до члена с номером $k$ и выглядит как $a_{1}, a_{2}, \dots, a_{k}$. Выберем наименьшее натуральное число $S$, которое еще не встречалось как разность между членами искомой последовательности, то есть $S \neq a_{i} - a_{j}$, где $j < i$ и $i, j \leqslant k$. Тогда добавим два новых элемента $a_{k+1} = 2a_{k}$ и $a_{k+2} = 2a_{k} + S$. Заметим, что $a_{k+1} - a_{i} = 2a_{k} - a_{i} \geqslant a_{k}$, где $i \leqslant k$, а значит точно не встречается среди разностей уже построенной последовательности. 
\end{Solution}

\chapter{1976}

\begin{exercise}[387]
Существует ли такое натуральное число, что если приписать его само к себе, то получится точный квадрат? 
\end{exercise}

\begin{Solution}
Обозначим искомое число $n$. Тогда задачу можно записать в виде
$$
n(10^{p}+1) = x^{2}
$$
И $n < 10^{p}$. Легко заметить, что если $10^{p} + 1$ содержит любой простой множитель более чем в первой степени, то искомое число нашлось. Докажем, что $10^{11}+1$ подходит. Действительно
$$
10^{11}+1 = 11 \cdot (10^{10} - 10^{9} + 10^{8} - 10^{7} + \dots + 1)
$$
Легко видеть, что выражение в скобках делится на 11. Значит искомое число
$$
n = \frac{10^{11} + 1}{121}
$$ 
\end{Solution}

\begin{exercise}[390]
Существует бесконечно много таких натуральных чисел $n$, что сумма цифр десятичной записи числа $2^{n}$ больше суммы цифр десятичной записи числа $2^{n+1}$. Докажите это.
\end{exercise}

\begin{exercise}[393]
Найдите сумму
$$
f(0) + f\left(\frac{1}{n}\right) + f\left(\frac{2}{n}\right) + \dots + f\left(\frac{n-1}{n}\right)+ f(1) 
$$
где $f(x) = \frac{4^{x}}{4^{x}+2}$.
\end{exercise}

\begin{Solution}
Сложим этот ряд с самим собой, но записанным в обратном порядке. Заметим что
$$
f\left(\frac{i}{n}\right) + f\left(\frac{n-i}{n}\right) = \frac{4^{\frac{i}{n}}}{4^{\frac{i}{n}}+2} + \frac{4^{\frac{n-i}{n}}}{4^{\frac{n-i}{n}}+2} = \frac{4 + 2 \cdot 4^{\frac{i}{n}} + 4 + 2 \cdot 4^{\frac{n-i}{n}}}{4 + 2 \cdot 4^{\frac{i}{n}} + 2 \cdot 4^{\frac{n-i}{n}} + 4} = 1
$$
То есть, удвоенная сумма ряда равна $n + 1$. Значит ответ $\frac{n+1}{2}$.
\end{Solution}

\chapter{1977}

\begin{exercise}[425]
Существует ли такое натуральное $n$, что каждое рациональное число между нулём и единицей представимо в виде суммы $n$ чисел, обратных натуральным?
\end{exercise}

\begin{exercise}[466]
Среди 1977 монет 50 фальшивых. Каждая фальшивая монета отличается от настоящей на один грамм (в ту или в другую сторону). Имеются чашечные весы со стрелкой, показывающей разность масс грузов на чашках. За одно взвешивание про одну выбранную монету нужно узнать, фальшивая она или настоящая. Научитесь это делать!
\end{exercise}

\chapter{1979}

\begin{exercise}[554]
Назовём натуральное число $n$ хорошим, если существуют (не обязательно различные) такие натуральные числа $a_{1}, a_{2}, \dots, a_{k}$, что $a_{1} + a_{2} + \dots + a_{k} = n$ и $\frac{1}{a_{1}} + \frac{1}{a_{2}} + \dots + \frac{1}{a_{k}} = 1$. Известно, что все числа между 33 и 73 - хорошие. Докажите, что все числа, большие 73, - тоже хорошие.
\end{exercise}

\begin{Solution}
Докажем несколько простых лемм. Во первых, если $n$ - хорошее, то $2n + 2$ тоже хорошее. 
$$
\sum_{i = 1}^{k} a_{i} = n \implies 2 + \sum_{i = 1}^{k} 2a_{i} = 2n + 2, \quad \sum_{i = 1}^{k} \frac{1}{a_{i}} = 1 \implies \frac{1}{2} + \sum_{i = 1}^{k} \frac{1}{2a_{i}} = 1
$$
Аналогично докажем, что $2n+9$ тоже хорошее
$$
\sum_{i = 1}^{k} a_{i} = n \implies 3 + 6 + \sum_{i = 1}^{k} 2a_{i} = 2n + 9, \quad \sum_{i = 1}^{k} \frac{1}{a_{i}} = 1 \implies \frac{1}{3} + \frac{1}{6} + \sum_{i = 1}^{k} \frac{1}{2a_{i}} = 1
$$
Далее очевидно по индукции.
\end{Solution}

\chapter{1980}

\begin{exercise}[644]
\begin{enumerate}
\item Существует выпуклый 1980-угольник со сторонами длины $1, 2, \dots, 1980$,
величины всех углов которого равны. Докажите это.
\item Существует ли такой 1981-угольник?
\end{enumerate}
\end{exercise}

\chapter{1984}

\begin{exercise}[882]
Если сумма трёх целых чисел равна нулю, то сумма их четвёртых степеней - удвоенный квадрат целого числа. Докажите это.
\end{exercise}

\begin{Solution}
Из $a + b + c = 0$ следует
$$
a^{2} + b^{2} + c^{2} = -2(ab + bc + ac) 
$$
А также
$$
(ab + bc + ac)^{2} = (ab)^{2} + (bc)^{2} + (ac)^{2}
$$
В итоге
$$
(a^{2} + b^{2} + c^{2})^{2} = a^{4} + b^{4} + c^{4} + 2(a^{2}b^{2} + b^{2}c^{2} + a^{2}c^{2}) \implies a^{4} + b^{4} + c^{4} = 2(ab + bc + ac)^{2}
$$
\end{Solution}

\chapter{1988}

\begin{exercise}[1088]
Если $pq + qr + pr = 1$, причём числа $p$, $q$ и $r$ рациональные, то число $(1 + p^{2})(1 + q^{2})(1 + r^{2})$ - квадрат рационального числа. Докажите это.
\end{exercise}

\begin{Solution}
Из условия следует
$$
1 + p^{2} = pq + qr + pr + p^{2} = (p + q)(p + r)
$$
Тогда, в силу симметрии
$$
(1 + p^{2})(1 + q^{2})(1 + r^{2}) = (q + r)^{2}(p + r)^{2}(p + q)^{2}
$$
\end{Solution}

\chapter{1990}

\begin{exercise}[1238]
Множество натуральных чисел разбито на два подмножества. В одном из них нет ни одной трёхчленной арифметической прогрессии. Обязательно ли в другом есть бесконечная арифметическая прогрессия?
\end{exercise}

\begin{exercise}[1244]
В сенате, состоящем из 30 сенаторов, каждые двое дружат или враждуют, причём каждый враждует ровно с 6 другими. Найдите общее количество троек сенаторов, в которых либо все трое дружат друг с другом, либо все трое враждуют между собой.
\end{exercise}

\begin{Solution}
Рассмотрим граф, где сенаторы это вершины, отношение дружбы это синее ребро, а отношение вражды это красное ребро. В этом графе $\binom{30}{3} = 4060$ треугольников. При этом если треугольник не одноцветный, то он имееет два красно-синих угла. А ребра при каждой вершине образуют $(29 - 6) \cdot 6 = 138$ красно-синих угла. Значит для количества $T$ не одноцветных треугольников можно записать соотношение $2 T = 30 \cdot 138$. Значит одноцветных треугольников $4060 - 138 \cdot 15 = 1990$.
\end{Solution}

\chapter{1991}

\begin{exercise}[1264]
На бесконечном листе белой клетчатой бумаги квадрат $2 \times 2$ нужно закрасить в чёрный цвет (а все остальные клетки должны остаться белыми). Можно ли это сделать несколькими операциями, каждая из которых - перекрашивание в противоположный цвет всех клеток квадрата $3 \times 3$ или $4 \times 4$?
\end{exercise}

\begin{exercise}[1275]
Последовательность $a_{1}, a_{2}, a_{3}, \dots$, натуральных чисел такова, что при любом натуральном $n$ верно равенство $a_{k+2} = a_{k} a_{k+1} + 1$. Докажите, что при $n > 9$ число $a_{n} - 22$ составное.
\end{exercise}

\begin{Solution}
Заметим, что если $a_{k} = 0$, то последовательность выглядит следующим образом
$$
a_{k-1}, 0, 1, 1, 2, 3, 7, 22, \dots 
$$
Но мы можем получить эту последовательность представляя, что это значения по модулю $a_{k}$. Значит 
$$
a_{k+6} = 22 \mod a_{k} \implies a_{k+6} - 22 = 0 \mod a_{k}
$$
\end{Solution}

\begin{exercise}[1280]
Период десятичного разложения дроби $\frac{1}{3^{100}}$ содержит
\begin{enumerate}
\item не менее 20 одинаковых цифр подряд;
\item последовательность цифр 123 456 789.
\end{enumerate}
Докажите это.
\end{exercise}

\begin{exercise}[1288]
Число $235^{2} + 972^{2} = 1000009$ - составное. Докажите это.
\end{exercise}

\begin{Solution}
Заметим, что $1000009 = 1000^{2} + 3^{2}$, то есть представимо в виде суммы квадратов двумя разными способами.
\end{Solution}

\chapter{1992}

\begin{exercise}[1321]
Ладья побывала во всех клетках доски размером $n \times n$ клеток. Докажите, что она изменила направление своего движения не менее $2n-2$ раз. (Ладья движется параллельно одной из сторон квадрата.)
\end{exercise}

\begin{exercise}[1338]
Укажите способ вычисления $2n$-го числа последовательности Фибоначчи не более чем за $6n$ операций сложения, умножения и вычитания.
\end{exercise}

\begin{exercise}[1357]
Числа $91! \cdot 1901! - 1$ и $92! \cdot 1900! + 1$ делятся на 1993. Докажите это.
\end{exercise}

\begin{Solution}

\end{Solution}

\chapter{1994}

\begin{exercise}[1422]
Числа 312 500 051 и 1 280 000 401 - составные. Докажите это.
\end{exercise}

\begin{Solution}
Обозначим $x = 50$. Тогда
$$
312 500 051 = x^{5} + x + 1 = (x^{3} - x^{2} + 1)(x^{2} + x + 1) = 122501 \cdot 2551
$$
Обозначим $x = 20$. Тогда
$$
1 280 000 401 = x^{7} + x^{2} + 1 = (x^{5} + x + 1)x^{2} - (x^{3} - 1) = (x^{2} + x + 1)(x^{5} - x^{4} + x^{2} - x + 1) = 421 \cdot 3040381
$$
\end{Solution}

\chapter{1995}

\begin{exercise}[1473]
Обозначим через $c_{n}$ первую цифру десятичной записи числа $2^{n}$. 
\begin{itemize}
\item Сколько единиц среди первых 1000 членов этой последовательности?
\item Докажите, что в последовательности $c_{1} = 2$, $c_{2} = 4$, $c_{3} = 8$, $c_{4} = 1$, $c_{5} = 3$, $\dots$ встречается ровно 57 различных подпоследовательностей $c_{k}c_{k+1}\dots c_{k+12}$ длины 13.
\end{itemize} 
\end{exercise}

\begin{exercise}[1499]
Число 2 представимо в виде суммы трёх четвёртых степеней рациональных чисел бесконечным множеством способов. Докажите это.
\end{exercise}

\begin{Solution}
Пусть $a + b = c$. Тогда легко показать, что
$$
a^{4} + b^{4} + c^{4} = 2(ac + bc - ab)^{2}
$$
Значит остается найти условия при которых $ac + bc - ab$ является полным квадратом некоторого натурального числа $D$. Пусть $a = u + v$, $b = u - v$. Тогда
$$
ac + bc - ab = (a + b)^{2} - ab = 4u^{2} - u^{2} + v^{2} = 3u^{2} + v^{2} = D^{2}
$$
Пусть $u = pq$. Тогда $(D - v)(D + v) = 3p^{2}q^{2}$. Это уравнение можно решить следующим образом:
$$
\begin{cases}
D + v = 3p^{2}\\
D - v = q^{2}
\end{cases} \implies D = \frac{3p^{2} + q^{2}}{2}, \quad v = \frac{3p^{2} - q^{2}}{2}, \quad a = pq + \frac{3p^{2} - q^{2}}{2}, \quad b = pq - \frac{3p^{2} - q^{2}}{2},\quad c = 2pq
$$
\end{Solution}

\begin{exercise}[1523]
Пусть $a_{1} = 1$, $a_{2} = 2$, $a_{3} = 3$, $a_{4} = 4$, $a_{5} = 5$ и $a_{n+1} = a_{1}a_{2} \dots a_{n} - 1$ при $n > 4$. Докажите равенство $a_{1}^{2} + a_{2}^{2} + a_{3}^{2} + \dots + a_{70}^{2} = a_{1} a_{2} a_{3} \dots a_{70}$.
\end{exercise}


\chapter{1998}

\begin{exercise}[1662]
Может ли десятичная запись куба натурального числа начинаться с цифр 1998?
\end{exercise}

\chapter{1999}

\begin{exercise}[1667]
Натуральный ряд разбит на два бесконечных множества чисел. Докажите, что сумма некоторых 100 чисел одного из этих множеств равна сумме некоторых 100 чисел другого множества.
\end{exercise}

\begin{exercise}[1703]
Если $a^{m} + b^{m} + c^{m} = 0$ и $a^{n} + b^{n} + c^{n} = 0$, где $m$ и $n$ - натуральные числа, то $abc = 0$. Докажите это.
\end{exercise}

\chapter{2000}

\begin{exercise}[1731]
В ряд нарисованы 60 звёздочек. Два игрока по очереди заменяют звёздочки на цифры. Какую именно из оставшихся звёздочек заменять на цифру, решает игрок, делающий очередной ход. Докажите, что второй игрок может играть так, чтобы полученное число (возможно, начинающееся на цифру 0) делилось на 13.
\end{exercise}

\begin{Solution}
Заметим, что любое число вида $\overline{abcabc} = 1001 \cdot \overline{abc}$, а значит делится на 13. Значит можно разбить зведочки на 10 блоков по 6 звездочек, а второй игрок может дублировать цифру написанную первым игроком, чтобы получился набор блоков указанного выше типа. 
\end{Solution}

\chapter{2001}

\begin{exercise}[1756]
Среди любых трёх из нескольких данных натуральных чисел можно выбрать два, одно из которых делится на другое. Докажите, что все числа можно так покрасить двумя красками, что из любых чисел одного цвета одно делится на другое.
\end{exercise}

\chapter{2002}

\begin{exercise}[1805]
В математической олимпиаде участвовали 21 мальчик и 21 девочка. Каждый
из них решил не более 6 задач. Для любого мальчика и для любой девочки
существует задача, которую решили и он, и она. Докажите, что существует задача, которую решили по крайней мере три мальчика и три девочки.
\end{exercise}

\begin{exercise}[1812]
Натуральные числа $a$, $b$ и $c$ таковы, что $\gcd(a^{2} - 1, b^{2} - 1, c^{2} - 1) = 1$. Докажите равенство 
$$
\gcd(ab + c, bc + a, ca + b) = \gcd(a, b, c)
$$
\end{exercise}

\begin{exercise}[1830]
В возрастающей последовательности натуральных чисел каждое число, начиная с 2002-го, является делителем суммы всех предыдущих чисел. Докажите, что в этой последовательности найдётся число, начиная с которого каждый член последовательности равен сумме всех предыдущих.
\end{exercise}

\begin{exercise}[1831]
В наборе 20 гирек, массы которых различны. Среди любых одиннадцати из них
можно найти две, сумма масс которых равна 100 граммам. Докажите, что сумма масс всех гирек равна 1 килограмму.
\end{exercise}

\begin{exercise}[1845]
Назовём несоседние натуральные числа $a$ и $b$ близкими, если $a^{2} - 1$ делится на $b$ и $b^{2} - 1$ делится на $a$.
\begin{enumerate}
\item Пусть $n > 1$. Докажите, что на отрезке $[n; 8n - 8]$ существует пара близких чисел.
\item Укажите такое $n > 1$, что на отрезке $[n; 8n - 9]$ нет ни одной пары близких чисел.
\end{enumerate} 
\end{exercise}

\chapter{2006}

\begin{exercise}[1981]
В клетках таблицы размером $11 \times 11$ расставлены натуральные числа от 1 до 121. Дима посчитал произведение чисел в каждой строке, а Саша — произведение чисел в каждом столбце. Могли ли они получить одинаковые наборы из 11 чисел?
\end{exercise}

\begin{Solution}
Рассмотрим простые числа из промежутка от 1 до 121, для которых в таблице нет кратных. Это 61, 67, 71, 73, 79, 83, 89, 97, 101, 103, 107, 109, 113. Их 13 штук. Значит в каком то столбце их не менее двух. Но если эти два числа встретились в одном столбце, то одно из них или его кратное должно быть и в одной строке с другим. Но этого быть не может.  
\end{Solution}

\begin{exercise}[2003]
Докажите следующие утверждения.
\begin{enumerate}
\item Для любых натуральных чисел $a$, $b$, $c$, $n$ уравнение $x^{2} + y^{2} + z^{2} = (a^{2} + b^{2} + c^{2})^{n}$ разрешимо в натуральных числах $x$, $y$, $z$.
\item Для любого нечётного числа $n \geqslant 3$ и любых натуральных чисел $a$, $b$, и $c$ уравнение $ax^{2} + by^{2} + cz^{2} = t^{n}$ разрешимо в натуральных числах $x$, $y$, $z$, $t$.
\item Существуют такие натуральные числа $a$, $b$ и $c$, что уравнение $ax^{2} + by^{2} + cz^{2} = t^{2}$ неразрешимо в натуральных числах $x$, $y$, $z$, $t$.
\end{enumerate}
\end{exercise}

\begin{exercise}[2023]
$a$, $b$, $c$ - ненулевые целые числа, сумма которых равна 0. Докажите, что
\begin{enumerate}
\item $(ab)^{5} + (bc)^{5} + (ca)^{5}$ делится на $(ab)^{2} + (bc)^{2} + (ca)^{2}$;
\item если $n - 1$ делится на 3, то $a^{n} + b^{n} + c^{n}$ делится на $a^{4} + b^{4} + c^{4}$;
\item если $n - 2$ делится на 3, то $(ab)^{n} + (bc)^{n} + (ca)^{n}$ делится на $(ab)^{2} + (bc)^{2} + (ca)^{2}$.
\end{enumerate}
\end{exercise}

\chapter{2007}

\begin{exercise}[2033]
У ведущего имеется колода из 52 карт. Зрители хотят узнать, в каком порядке лежат карты (не уточняя, сверху вниз или снизу вверх). Разрешено задавать ведущему вопросы вида "Сколько карт лежит между такой-то и такой-то картами?". Один из зрителей знает, в каком порядке лежат карты. Какое наименьшее число вопросов он должен задать, чтобы остальные зрители по ответам на эти вопросы могли узнать порядок карт в колоде?
\end{exercise}

\chapter{2008}

\begin{exercise}[2096]
Депутаты парламента образовали 2008 комиссий, каждая - не более чем из 10 человек. Любые 11 комиссий имеют хотя бы одного общего члена. Докажите, что существует человек, входящий во все комиссии.
\end{exercise}

\begin{exercise}[2097]
Найдите все такие простые числа $p$ вида $a^{2} + b^{2} + c^{2}$, что $a^{4} + b^{4} + c^{4}$ делится на $p$.
\end{exercise}

\begin{exercise}[2114]
Существует бесконечно много таких натуральных чисел, не делящихся на 10, что все ненулевые цифры десятичных записей их квадратов нечётны. Докажите это.
\end{exercise}


\chapter{2010}

\begin{exercise}[2162]
Семизначный код, состоящий из семи различных цифр, назовём хорошим. Паролем сейфа является хороший код. Сейф откроется, если введён хороший код и на каком-нибудь месте цифра кода совпала с соответствующей цифрой пароля. За какое наименьшее количество попыток можно с гарантией открыть сейф?
\end{exercise}

\begin{exercise}[2202]
Дана арифметическая прогрессия из 22 различных натуральных чисел, каждое из которых является точной степенью (то есть степенью натурального числа, большей 1). Докажите, что разность этой прогрессии больше 2010.
\end{exercise}

\begin{Solution}
Рассмотрим прогрессию $a_{1}, a_{2}, a_{3}, \dots, a_{22}$ с разностью $d$. Докажем, что $d$ делится на каждое простое из множества $S = \{2, 3, 5, 7, 11\}$. Пусть это не так и $d$ не делится на некоторое $q \in S$. По принципу Дирихле, среди первых $q$ членов прогрессии найдется $a_{n}$, делящийся на $q$. Но тогда для $m = n + q \leqslant 2q = 22$,  $a_{m}$ тоже делится на $q$. Тогда, по условию, $a_{n}$ и $a_{m}$ делятся на $q^{2}$. А значит и разность между ними тоже делится на $q^{2}$. Но это противоречит тому, что $d$ не делится на $q$. Значит $d$ делится на $2 \cdot 3 \cdot 5 \cdot 7 \cdot 11 = 2310$.
\end{Solution}

\chapter{2020}



\begin{exercise}[2593]
Каждая вершина правильного многоугольника окрашена в один из трех цветов так, что в каждый из трех цветов окрашено нечетное число вершин. Докажите, что количество равнобедренных треугольников, вершины которых окрашены в три разных цвета, нечетно.
\end{exercise}

\begin{exercise}[2601]
Глеб задумал натуральные числа $N$ и $a$, где $a < N$. Число $a$ он написал на доске. Затем Глеб стал проделывать такую операцию: делить $N$ с остатком на последнее выписанное на доску число и полученный остаток от деления также записывать на доску. Когда на доске
появилось число 0, он остановился. Мог ли Глеб изначально выбрать такие $N$ и $a$, чтобы сумма выписанных на доске чисел была больше $100N$?
\end{exercise}

\chapter{2021}

\begin{exercise}[2638]
Существует ли целое положительное число $n$ такое, что все его цифры (в десятичной записи) больше 5, а все цифры числа $n^{2}$ меньше 5?
\end{exercise}

\begin{exercise}[2669]
Докажите, что для любого натурального числа $n$ числа $1, 2, \dots, n$ можно разбить на несколько групп так, чтобы сумма чисел в каждой группе была равна степени тройки.
\end{exercise}

\begin{exercise}[2673]
В очереди на посадку в $n$-местный самолет стоят $n$ пассажиров. Первой в очереди стоит рассеянная старушка, которая, зайдя в самолет, садится на случайно выбранное место. Каждый следующий пассажир садится на свое место, если оно свободно, и на случайное место в противном случае. Сколько в среднем пассажиров окажутся не на своих местах?
\end{exercise}

\chapter{2024}

\begin{exercise}[2810]
Дано натуральное число $n \geqslant 2$. Сколькими способами можно раскрасить клетки квадрата $n \times n$ в четыре цвета так, чтобы любые две клетки, имеющие общую сторону или вершину, были покрашены в
разные цвета?
\end{exercise}


\end{document} 
 
 
